\section{Outputs}\label{sec:outputs}

What does a run leave, and how is it read? Five tables in Parquet (Table~\ref{tab:outputs});
the main one, T3, cuts each flight into stretches of constant membership.

\begin{table}
\centering
\small
\begin{tabularx}{\textwidth}{@{}l X@{}}
\toprule
Table & Columns (type) and content \\
\midrule
T1 \texttt{pair\_intervals} & \texttt{date}, \texttt{a}, \texttt{b} (int32), \texttt{t\_on}, \texttt{t\_off} (int32, s),
  \texttt{duration} (int32), \texttt{ell\_med} (int16, s, signed: positive if \texttt{b} follows \texttt{a}),
  \texttt{d\_dec} (float32, m). Who was with whom, and when. \\ \addlinespace
T2 \texttt{configurations} & \texttt{date}, \texttt{conf\_id}, \texttt{episode\_id} (int32), \texttt{t\_start},
  \texttt{t\_end} (int32), \texttt{n} (int16), \texttt{members} (list of int32). \\ \addlinespace
T3 \texttt{membership} & \texttt{flight} (int32), \texttt{t\_start}, \texttt{t\_end} (int32), \texttt{state} (int8:
  L, S, G), \texttt{conf\_id}, \texttt{episode\_id} (int32, null if L or S), \texttt{n} (int16, 1 if S).
  The main table: each flight cut into stretches of constant membership. \\ \addlinespace
T4 \texttt{episodes} & \texttt{episode\_id}, \texttt{date}, \texttt{t\_start}, \texttt{t\_end}, \texttt{n\_flights},
  \texttt{n\_max}, \texttt{n\_config}, \texttt{block}. \\ \addlinespace
T5 \texttt{flight\_summary} & \texttt{flight}, \texttt{phi\_G} (fraction of $T^\dep_i$ in a group), \texttt{kappa}
  (distinct companions), \texttt{n\_entries}, \texttt{n\_episodes}, \texttt{block}, \texttt{local} (boolean). \\
\bottomrule
\end{tabularx}
\caption{Final tables, in Parquet. The members of a configuration are not repeated in T3: they are
read from T2 through \texttt{conf\_id}.}
\label{tab:outputs}
\end{table}

The state of flight $i$ at instant $t$ is obtained with an as-of join on (\texttt{flight}, $t$)
in T3, sorted by \texttt{t\_start} (Figure~\ref{fig:tables}). T3 has about
$N_{\rm flights}(2+2\bar c)$ rows, with $\bar c$ the mean number of configurations per flight:
of the order of $10^5$--$10^6$ rows, tens of MB. The other tables are smaller.

\begin{figure}
\centering
\begin{tikzpicture}[x=1cm,y=1cm,
    tbl/.style={draw=ink2, fill=white, rounded corners=1.5pt, align=left, font=\footnotesize, inner sep=4pt, anchor=north west},
    key/.style={->, blueI, line width=0.8pt}]
  \node[tbl] (t1) at (0,3.6) {\textbf{T1} pair\_intervals\\ a, b, t\_on, t\_off};
  \node[tbl] (t2) at (4.6,3.6) {\textbf{T2} configurations\\ \underline{conf\_id}, episode\_id,\\ t\_start, t\_end, members};
  \node[tbl] (t3) at (4.6,0.8) {\textbf{T3} membership\\ flight, t\_start, t\_end,\\ state, conf\_id};
  \node[tbl] (t4) at (10.6,3.6) {\textbf{T4} episodes\\ \underline{episode\_id}, block};
  \node[tbl] (t5) at (10.6,0.8) {\textbf{T5} flight\_summary\\ flight, phi\_G, block};
  \draw[->, ink2, line width=0.8pt] (t1.east) ++(0,-0.15) -- node[above, note] {Phase 4} ($(t2.west)+(0,-0.15)$);
  \draw[key] (t3.north) -- node[right, note] {conf\_id} (t2.south);
  \draw[key] (t2.east) ++(0,-0.15) -- node[above, note] {episode\_id} ($(t4.west)+(0,-0.15)$);
  \draw[key] (t5.north) -- node[right, note] {block} (t4.south);
  \node[note, anchor=north west, align=left] at (0,0.8) {state of flight 8 at 13:42:\\the row of T3 with\\t\_start $\le$ 13:42 $<$ t\_end,\\then its members in T2};
\end{tikzpicture}
\caption{How the tables connect. Arrows point from a column to the table where it is a key. The
example reads the state of a flight at one instant with an as-of join on T3.}
\label{fig:tables}
\end{figure}
